\begin{abstract}
Raport opisuje pamięć skojarzeniową opartą na sieci Hopfielda, w której porównano dwie metody uczenia: regułę Hebba i regułę Oji w wariancie podprzestrzeni. Przedstawiono model sieci i obie reguły razem z wyjaśnieniem, dlaczego działają, oraz program złożony z warstwy obliczeniowej w języku Rust i warstwy wizualizacji w języku Python. Program symuluje sieć w trybie synchronicznym i asynchronicznym i pokazuje kolejne stany jako bitmapy. Na dostarczonych zbiorach bitmap zbadano stabilność wzorców uczących, odzyskiwanie wzorców z wejścia z szumem i z zasłonięciem oraz reakcję sieci na losowe wejście. Skonstruowano też liczny zbiór bitmap \mbox{5x5}, w którym wszystkie wzorce są stabilne, i przykład oscylacji między dwoma stanami. Na koniec sprawdzono działanie sieci na zdjęciach kotów o rozmiarze \mbox{300x200} pikseli. Reguła Hebba zawodzi na zbiorach skorelowanych wzorców, nawet gdy wzorców jest mało, i stabilizuje tylko 12 z 74 wzorców dostarczonych zbiorów. Reguła Oji zbiega do rzutu na podprzestrzeń wzorców, dlatego stabilizuje je wszystkie i odzyskuje zdjęcia kotów także z silnie zaszumionego wejścia.
\end{abstract}
